% Circuit QED experimental proposal
\documentclass[aps,prl,twocolumn,superscriptaddress,nofootinbib]{revtex4-2}

\usepackage{amsmath,amssymb}
\usepackage{hyperref}
\usepackage{booktabs}

\begin{document}

\title{Lamb Shift Dependence on Resonator Mode Structure:\\
$\lambda/4$ versus $\lambda/2$ in Circuit QED}

\author{Wright Brothers Research Program}
\affiliation{Independent Research, 2026}

\date{\today}

\begin{abstract}
We propose a direct comparison of the Lamb shift experienced by
a transmon qubit coupled to a $\lambda/4$ (quarter-wave) versus
a $\lambda/2$ (half-wave) superconducting coplanar or coaxial
resonator.  The $\lambda/2$ resonator supports all harmonics
($f, 2f, 3f, \ldots$), while the $\lambda/4$ supports only odd
harmonics ($f, 3f, 5f, \ldots$).  The vacuum fluctuations of the
``missing'' even modes contribute to the qubit Lamb shift in the
$\lambda/2$ case but not in the $\lambda/4$ case.  We compute the
expected difference $\Delta(\delta f)\sim 1$\,MHz for typical
circuit QED parameters and show it is well above the measurement
floor ($\sim 1$\,Hz with standard techniques).  This measurement
provides a direct quantification of multi-mode vacuum fluctuation
contributions to qubit frequency, relevant to understanding
frequency stability and residual shifts in multi-mode quantum
processors.  No new hardware is required: both resonator types
and transmon qubits are standard components in superconducting
circuit laboratories.
\end{abstract}

\maketitle

%===================================================================
\section{Motivation}
%===================================================================

Superconducting transmon qubits are routinely coupled to
microwave resonators in the dispersive
regime~\cite{Blais2004,Koch2007}.  The qubit frequency acquires
a shift from virtual photon exchange with each resonator mode
(the Lamb shift):
\begin{equation}\label{eq:lamb}
  \delta f = \sum_n \frac{g_n^2}{f_q - f_n},
\end{equation}
where $g_n$ is the coupling to mode $n$, $f_q$ the bare qubit
frequency, and $f_n$ the $n$-th resonator mode
frequency~\cite{Schuster2007,Gambetta2006}.

In multi-mode environments---a regime of growing importance for
quantum processors with bus resonators, multiplexed readout, and
3D cavities~\cite{Paik2011,Reagor2016}---the sum in
Eq.~\eqref{eq:lamb} extends over many modes, and higher modes
make non-negligible contributions.  Accurately predicting and
controlling these multi-mode shifts is essential for qubit
frequency targeting and cross-talk minimization.

Despite this practical importance, \textbf{no experiment has
directly compared the Lamb shift in resonators with different
mode structures}.  We propose the simplest such comparison:
$\lambda/4$ versus $\lambda/2$.

%===================================================================
\section{Mode Structure}
%===================================================================

A $\lambda/2$ resonator (both ends shorted, or both open) supports
all integer harmonics:
\begin{equation}
  f_n^{(\lambda/2)} = n\,f_0,\quad n=1,2,3,4,\ldots
\end{equation}
A $\lambda/4$ resonator (one end shorted, one end open) supports
only odd harmonics:
\begin{equation}
  f_n^{(\lambda/4)} = (2n{-}1)\,f_0,\quad n=1,2,3,\ldots
\end{equation}
Both have the same fundamental frequency $f_0$ if the $\lambda/4$
resonator is half the physical length of the $\lambda/2$ resonator
(accounting for the factor of 2 in the standing-wave condition).

The two resonators share all odd-harmonic modes; they differ only
in the presence (absence) of even harmonics in the $\lambda/2$
($\lambda/4$) case.

%===================================================================
\section{Predicted Lamb Shift Difference}
%===================================================================

The Lamb shift difference between the two configurations isolates
the even-mode vacuum contribution:
\begin{equation}\label{eq:diff}
  \Delta(\delta f) = \delta f^{(\lambda/2)} - \delta f^{(\lambda/4)}
  = \sum_{n\;\mathrm{even}} \frac{g_n^2}{f_q - f_n}.
\end{equation}

For typical circuit QED parameters ($f_0 = 6\;\mathrm{GHz}$,
$f_q = 5\;\mathrm{GHz}$, $g/2\pi = 100\;\mathrm{MHz}$), we
estimate the individual even-mode contributions (Table~\ref{tab:modes}).

\begin{table}[h]
\caption{Even-mode Lamb shift contributions
for $f_0=6\;\mathrm{GHz}$, $f_q=5\;\mathrm{GHz}$,
$g_n\approx g_0\sqrt{f_0/f_n}$.}
\label{tab:modes}
\begin{ruledtabular}
\begin{tabular}{cccc}
$n$ & $f_n$ (GHz) & $g_n/2\pi$ (MHz) & $g_n^2/(f_q{-}f_n)$ (kHz) \\
\hline
2 & 12 & 71 & $-714$ \\
4 & 24 & 50 & $-132$ \\
6 & 36 & 41 & $-54$ \\
8 & 48 & 35 & $-29$ \\
\hline
\multicolumn{3}{l}{Total $\Delta(\delta f)$} & $\approx -930$\,kHz \\
\end{tabular}
\end{ruledtabular}
\end{table}

The predicted difference $|\Delta(\delta f)|\approx 930\;\mathrm{kHz}
\approx 1\;\mathrm{MHz}$ is large compared to the spectroscopic
resolution achievable with Ramsey or echo techniques
($\sim 1\;\mathrm{Hz}$ with one hour of
averaging~\cite{Sank2016}).  The signal-to-noise ratio exceeds
$10^5$.

%===================================================================
\section{Experimental Protocol}
%===================================================================

\begin{enumerate}
\item \textbf{Sample preparation.}  Fabricate (or select from
existing stock) a $\lambda/2$ and a $\lambda/4$ superconducting
resonator with matched fundamental frequencies.  Standard
aluminum coplanar waveguide or coaxial stub resonators suffice.
Couple each to a transmon qubit with $g/2\pi\sim 100\;\mathrm{MHz}$.

\item \textbf{Cooldown.}  Mount in a dilution refrigerator at
$T\lesssim 20\;\mathrm{mK}$.

\item \textbf{Characterization.}  Measure $S_{21}(f)$ to verify
mode frequencies:
\begin{itemize}
\item $\lambda/2$: peaks at $f_0, 2f_0, 3f_0, \ldots$
\item $\lambda/4$: peaks at $f_0, 3f_0, 5f_0, \ldots$ (even modes absent)
\end{itemize}

\item \textbf{Qubit spectroscopy.}  Measure $f_q$ in each
configuration via two-tone spectroscopy or Ramsey interferometry.

\item \textbf{Comparison.}  Extract
$\Delta(\delta f) = f_q^{(\lambda/2)} - f_q^{(\lambda/4)}$
and compare with the theoretical prediction from
Eq.~\eqref{eq:diff}.
\end{enumerate}

If the same physical qubit can be coupled to both resonators
sequentially (e.g., via a tunable coupler or by physically
remounting), systematic errors from qubit-to-qubit variation
are eliminated.

%===================================================================
\section{What This Measurement Tests}
%===================================================================

\subsection{Multi-mode vacuum contribution}

The measurement directly quantifies how much of the qubit's Lamb
shift comes from higher (even) harmonics.  This is relevant for:
\begin{itemize}
\item \textbf{Frequency targeting}: predicting the dressed qubit
frequency in multi-mode environments.
\item \textbf{Cross-talk}: understanding mode-mediated interactions
between qubits in bus architectures.
\item \textbf{Decoherence}: multi-mode Purcell decay rates depend
on the same couplings $g_n$ that produce the Lamb shift.
\end{itemize}

\subsection{Completeness of the mode-sum description}

Standard circuit QED predicts Eq.~\eqref{eq:diff} as a simple sum
of individual even-mode contributions.  If the measured
$\Delta(\delta f)$ deviates from this prediction, it would indicate
a collective or non-perturbative vacuum effect beyond the
single-mode-sum picture.

Such a deviation is not expected by standard theory, but
\emph{has never been explicitly tested}.  This measurement provides
the first direct comparison of vacuum fluctuation effects in
resonators with structurally different mode spectra.

%===================================================================
\section{Relation to Prior Work}
%===================================================================

Lamb shifts and dispersive shifts in circuit QED are well
studied~\cite{Schuster2007,Gambetta2006,Paik2011}, and
multi-mode effects have been analyzed
theoretically~\cite{Nigg2012,Solgun2014}.  Measurements in
3D cavities~\cite{Reagor2016} and multi-mode
environments~\cite{McKay2015} have characterized individual mode
couplings.

The novelty of the present proposal is the \emph{controlled
comparison} between two resonators differing only in mode
parity---one containing all harmonics, the other only odd
harmonics---providing a clean test of the vacuum mode sum.

We note that the mode structure of a $\lambda/4$ resonator
(odd harmonics only: $1, 3, 5, \ldots$) is equivalent to the
mode structure of a cavity with Dirichlet--Neumann (DN) mixed
boundary conditions, whose vacuum energy has the opposite sign
to the standard Dirichlet--Dirichlet (DD) case---a result
first obtained by Boyer~\cite{Boyer1974}.  The present
measurement can thus be viewed as a circuit QED realization
of Boyer's fifty-year-old prediction, probing whether the
different mode-parity structure leads to measurably different
vacuum physics.

%===================================================================
\section{Conclusion}
%===================================================================

We propose comparing the Lamb shift of a transmon qubit in
$\lambda/4$ and $\lambda/2$ superconducting resonators.  The
predicted difference ($\sim 1$\,MHz) is five orders of magnitude
above the measurement floor.  The experiment requires no new
hardware and can be performed with standard components available
in any superconducting circuit laboratory.  A positive result
validates the multi-mode vacuum contribution to qubit frequency;
any deviation from the single-mode-sum prediction would reveal
collective vacuum effects not captured by standard perturbation
theory.

\begin{thebibliography}{99}
\bibitem{Blais2004} A.~Blais, R.-S.~Huang, A.~Wallraff, S.~M.~Girvin, and R.~J.~Schoelkopf, Phys.\ Rev.\ A \textbf{69}, 062320 (2004).
\bibitem{Koch2007} J.~Koch \textit{et al.}, Phys.\ Rev.\ A \textbf{76}, 042319 (2007).
\bibitem{Schuster2007} D.~I.~Schuster \textit{et al.}, Nature \textbf{445}, 515 (2007).
\bibitem{Gambetta2006} J.~Gambetta \textit{et al.}, Phys.\ Rev.\ A \textbf{74}, 042318 (2006).
\bibitem{Paik2011} H.~Paik \textit{et al.}, Phys.\ Rev.\ Lett.\ \textbf{107}, 240501 (2011).
\bibitem{Reagor2016} M.~Reagor \textit{et al.}, Phys.\ Rev.\ B \textbf{94}, 014506 (2016).
\bibitem{Sank2016} D.~Sank \textit{et al.}, Phys.\ Rev.\ Lett.\ \textbf{117}, 190503 (2016).
\bibitem{Nigg2012} S.~E.~Nigg \textit{et al.}, Phys.\ Rev.\ Lett.\ \textbf{108}, 240502 (2012).
\bibitem{Solgun2014} F.~Solgun, D.~W.~Abraham, and D.~P.~DiVincenzo, Phys.\ Rev.\ B \textbf{90}, 134504 (2014).
\bibitem{McKay2015} D.~C.~McKay \textit{et al.}, Phys.\ Rev.\ Applied \textbf{8}, 034004 (2017).
\bibitem{Boyer1974} T.~H.~Boyer, Phys.\ Rev.\ A \textbf{9}, 2078 (1974).
\end{thebibliography}

\end{document}
